
\section{\trGenus{} $4$ 与 $5$ 的\trCurve}

我们先考虑\trGenus{} $4$ 的非超椭圆\trCurve{}$C$。注意 $\deg\omega_C=6$ 且 $h^0(C,\omega_C)=4$，所以\trCanonical{}映射把 $C$ 表为 $\mathbb{P}^3$ 中的一条六次曲线。设 $\iota\colon C\hookrightarrow\mathbb{P}^3$ 为\trCanonical{}（闭）嵌入。我们将会看到：所有这类 $C$ 都是\trQuadricSurface{}与三次曲面的\trCompleteIntersection；反之，二次曲面与三次曲面的一切\trRegular{}\trCompleteIntersection{}都是\trGenus{} $4$ 的非超椭圆\trCurve，且以\trCanonical{}方式嵌入。

由 \eqref{eq:h0-high-degree}（\trRiemannRoch{}与\trSerreDuality），
\[
h^0(C,\iota^*\OO_{\mathbb{P}^3}(2))=h^0(C,\omega_C^{\otimes 2})=\deg\omega_C^{\otimes 2}-g+1=12-4+1=9.
\]
另一方面，$h^0(\mathbb{P}^3,\OO_{\mathbb{P}^3}(2))=\binom{3+2}{2}=10$（\mbox{练习~15.1.C}）。于是\trRestrictionMap
\[
\mathrm{H}^0(\mathbb{P}^3,\OO_{\mathbb{P}^3}(2))\longrightarrow \mathrm{H}^0(C,\iota^*\OO_{\mathbb{P}^3}(2))
\]
必有非平凡\trKernel，故 $\mathbb{P}^3$ 中至少存在一张包含我们\trCurve{}$C$ 的\trQuadricSurface{} $Q$。\trQuadricSurface{}要么是\trDoublePlane，要么是两个平面之并，要么是锥，要么是\trRegular{}\trQuadricSurface。（它们分别对应\trRank{}为 $1$、$2$、$3$、$4$ 的\trQuadraticForm{}；\trQuadraticForm{}之\trRank{}的定义见\mbox{练习~5.4.J}。）但 $C$ 不可能落在一张平面内，所以 $Q$ 必是锥或\trRegular{}\trQuadricSurface。特别地，$Q$ \trIrreducible。

现在，$C$ 不可能落在两张（互不相同的）此类二次曲面上，设为 $Q$ 与 $Q'$。否则，由于 $Q$ 与 $Q'$ 没有公共\trIrreducibleComponent（它们\trIrreducible{}且互不相同！），$Q\cap Q'$ 是一条\trCurve（不必\trReduced{}或\trIrreducible）。由\trBezout{}定理（\mbox{练习~18.6.J}），$Q\cap Q'$ 是一条\trDegree{}为 $4$ 的\trCurve。于是我们的\trCurve{}$C$ 既然\trDegree{}为 $6$，就不可能包含在 $Q\cap Q'$ 之中。（若你不能直接看出这一点，\mbox{练习~18.6.E} 或许有帮助。）

\textbf{19.8.1.}\quad 接下来我们考虑三次曲面。再由 \eqref{eq:h0-high-degree}，
\[
h^0(C,\omega_C^{\otimes 3})=\deg\omega_C^{\otimes 3}-g+1=18-4+1=15.
\]
而 $\Sym^3\Gamma(C,\omega_C)$ 的维数为 $\binom{3+3}{3}=20$。于是 $C$ 至少落在一个五维的三次形式向量空间之中。其中四维的子空间来自用一个线性形式 $(?w+?x+?y+?z)$\footnote{译注：原书以问号表示未定系数。}去乘\trQuadraticForm{} $Q$。故仍剩下一个三次曲面 $K$，其底层形式不被\trQuadraticForm{} $Q$ 整除\footnote{原书此处作「quadric form」，疑当为「quadratic form」之误，兹照录原书。——译者注}（即 $K$ 不包含 $Q$）。于是 $K$ 与 $Q$ 无公共\trIrreducibleComponent，故 $K\cap Q$ 是一个包含 $C$ 作为\trClosedSubscheme{}的\trCompleteIntersection。$K\cap Q$ 与 $C$ 的\trDegree{}都是 $6$（前者由\trBezout{}定理，\mbox{练习~18.6.J}；后者因为 $C$ 由一个\trDegree{}为 $6$ 的\trLineBundle{}嵌入，\mbox{练习~18.6.H}）。此外，$K\cap Q$ 与 $C$ 的\trArithmeticGenus{}都是 $4$（前者由\mbox{练习~18.6.S}）。这两个不变量决定（线性）\trHilbert{}多项式，故 $K\cap Q$ 与 $C$ 有相同的\trHilbert{}多项式。于是由\mbox{练习~18.6.E}，$C=K\cap Q$。

我们现在证明逆命题：任一个二次曲面与一个三次曲面之\trCompleteIntersection{}$C$，只要是\trRegular{}的，就是一条以\trCanonical{}方式嵌入的\trGenus{} $4$ \trCurve。由\mbox{练习~18.6.S}，这样的\trCompleteIntersection{}的\trGenus{}为 $4$。

\textbf{19.8.A.（练习）}\quad 证明：$\OO_C(1)$ 至少有四个截面。（换个说法：$C$ 不落在任何\trHyperplane{}内。）

\trDegree{}为 $2g-2$ 且带有（至少）$g$ 个截面的\trInvertibleSheaf{}只有\trCanonical{}层（\mbox{练习~19.2.A}），所以 $\OO_C(1)\cong\omega_C$，$C$ 确实是以\trCanonical{}方式嵌入的。

\textbf{19.8.B.（练习）}\quad 给出一个启发式论证，表明非超椭圆的\trGenus{} $4$ \trCurve{}「构成一个维数为 $9$ 的族」。

现在向\trGenus{} $5$ 进发！

\textbf{19.8.C.（练习）}\quad 设 $C$ 是一条非超椭圆的\trGenus{} $5$ \trCurve。证明：\trCanonical{}曲线是 $\mathbb{P}^4$ 中\trDegree{}为 $8$ 的曲线。证明：它落在某个三维的\trQuadricSurface{}向量空间中（即它落在三张线性无关的\trQuadricSurface{}上）。证明：三张\trQuadricSurface{}的\trRegular{}\trCompleteIntersection{}是一条\trCanonical{}（嵌入）的\trGenus{} $5$ \trCurve。

遗憾的是，并非所有\trCanonical{}的\trGenus{} $5$ 曲线都是 $\mathbb{P}^4$ 中三张\trQuadricSurface{}的\trCompleteIntersection。但在「大多数\trGenus{} $3$ \trCurve{}可以描述为平面四次曲线」的同一意义上，大多数\trCanonical{}的\trGenus{} $5$ 曲线是三张\trQuadricSurface{}的\trCompleteIntersection，而大多数\trGenus{} $5$ \trCurve{}是非超椭圆的。说这件事的正确方式是：在\trGenus{} $5$ \trCurve{}的\trModuliSpace{}中存在一个稠密的\trZariski{}开轨迹，由\trCanonical{}嵌入被三张\trQuadricSurface{}切出的非超椭圆\trCurve{}组成。

（那些不被一个三维\trQuadricSurface{}向量空间切出的非超椭圆\trGenus{} $5$ 典范曲线，恰是三角\trCurve；见\mbox{练习~19.7.C}。在三角映射下被映到 $\mathbb{P}^1$ 同一点的三点组，结果落在\trCanonical{}嵌入中的一条直线上。凡沿这三点消失的\trQuadricSurface{}必沿该直线消失——说白了，任何带三个零点的二次多项式必是零多项式。）

\textbf{19.8.D.（练习）}\quad 假定上面的讨论，数一数三张\trQuadricSurface{}的\trCompleteIntersection，给出一个启发式论证，表明\trGenus{} $5$ \trCurve{}「构成一个维数为 $12$ 的族」。

\textbf{19.8.2.}\quad 通过把\trCanonical{}曲线当作\trCompleteIntersection{}来思考，我们已经理解了\trGenus{} $3$ 到 $5$ 的\trCurve。遗憾的是，我们的好运气到头了。

\textbf{19.8.E.（练习）}\quad 证明：若 $C\subset\mathbb{P}^{g-1}$ 是\trGenus{} $g\ge 6$ 的\trCanonical{}曲线，则 $C$ \emph{不是}\trCompleteIntersection。（提示：\trBezout{}定理，\mbox{练习~18.6.J}。）

\textbf{19.8.3.}\quad 依然有美妙的事情在上演。例如，对 $k=\bar{k}$ 上大多数\trGenus{} $6$ \trCurve{}而言，\trCanonical{}嵌入（$\mathbb{P}^5$ 中\trDegree{} $10$）按如下方式产生。从 $G(2,5)$（维数 $6$；你能看出为什么吗？）在 $\mathbb{P}^9$ 中经其\trPluckerEmbedding{}之下的\trCompleteIntersection{}出发。将它与一个\trCodimension{}为 $4$ 的\trLinearSpace{}相交，得到 $\mathbb{P}^5$ 中一张\trSmooth{}曲面。这张曲面\trIsomorphic{}于在四个点处\trBlowup{}的 $\mathbb{P}^2$。（它是一张 \textbf{\trDelPezzo{}曲面}（del Pezzo surface），将在\mbox{练习~21.5.H} 中定义。这张特定的曲面出现在\mbox{注记~27.3.7} 的表中。）用一个\trQuadricHypersurface{}去切这张曲面（即与之相交），便得到一条\trGenus{} $6$ \trCurve。

\textbf{19.8.4.}\ \emph{关于一般\trGenus{}\trCurve{}的一些讨论。}\quad 然而，我们手头还是有一些数据的。若 $\mathscr{M}_g$ 是这个尚未良定义的「\trGenus{} $g$ \trCurve{}的\trModuliSpace」，我们有一些求其在低 $g$ 时维数的启发式方法。\trGenus{} $0$ 时，在\trAlgClosedField{}上只有一条\trGenus{} $0$ \trCurve（\mbox{命题~19.3.1}），于是看起来 $\dim\mathscr{M}_0=0$。\trGenus{} $1$ 时，在\trAlgClosedField{}上，我们很快就会看到\trEllipticCurve{}由 $j$-不变量分类（\mbox{练习~19.10.G}），于是看起来 $\dim\mathscr{M}_1=1$。我们还非正式地算得 $\dim\mathscr{M}_2=3$、$\dim\mathscr{M}_3=6$、$\dim\mathscr{M}_4=9$、$\dim\mathscr{M}_5=12$。规律是什么？事实上，在某种强意义上，\trRiemann{}早已知道：当 $g>1$ 时 $\dim\mathscr{M}_g=3g-3$。\trGenus{} $0$ 与\trGenus{} $1$ 时哪里出了问题？作为一条线索，回想讨论\trHilbert{}函数（\S18.6）时我们的洞见：每当某个函数「最终是多项式」，我们就该假定它「想成为多项式」，并且存在某个更好的函数（通常是\trEulerCharacteristic）\emph{是}多项式，而上同调消失保证原函数与那个更好的函数「最终一致」。在 $\mathscr{M}_g$ 的情形把这件事说清楚，远远超出我们当前讨论的范围，所以我们满足于观察如下事实。\emph{每一条}\trGenus{}大于 $1$ 的\trRegular{}\trCurve{}都有有限多个\trAutomorphism——一个零维的\trAutomorphismGroup。\emph{每一条}\trGenus{} $1$ 的\trRegular{}\trCurve{}都有一维的\trAutomorphismGroup（见\mbox{问题~19.10.6}）。而唯一那条\trGenus{} $0$ 的\trRegular{}\trCurve{}具有三维的\trAutomorphismGroup（\mbox{练习~15.5.B}）。（更多讨论见\mbox{旁注~21.5.12}。）于是注意到：对所有 $g\ge 0$，
\[
\dim\mathscr{M}_g-\dim\Aut C_g=3g-3,
\]
其中 $\Aut C_g$ 指任何一条\trGenus{} $g$ \trCurve{}的\trAutomorphismGroup。

事实上，用\trStack（或\trOrbifold）的语言来说，（光滑、射影、\trGeometricallyIrreducible{}的）\trGenus{} $0$ \trCurve{}的\trModuliSpace{}维数为 $-3$，而\trGenus{} $1$ \trCurve{}的\trModuliSpace{}维数为 $0$，这样说是有意义的。



